\subsection{Direct dense-kernel optimization controls}

Compression fidelity and optimization quality answer different questions.
The six-dataset controls compare the factorized and direct dense
implementations with identical starting assignments, capacity constraints,
search routines and iteration limits. All final quality values are evaluated
with the original kernel. At rank 256 and a 1\% size corridor, the mean ratio
of the factor solution's true cut to the dense solution's true cut is
\ControlsCutRatio{}, averaging seeds within each dataset before averaging the
six datasets; per-dataset quality and timing values are in the supplement.
Neither heuristic solution establishes an optimum. Higher rank
can improve objective fidelity without producing a monotonically better
optimized partition.

The comparison also exposes a size-dependent computational benefit. Dense
optimization is faster on the smaller molecular datasets, while compact
factors become faster as the number of entities grows. At 32,000 ChEMBL
fingerprints, rank 256 takes \ControlsLargeFactorSeconds{} seconds versus
\ControlsLargeDenseSeconds{} seconds for the dense arm, with mean peak process
memory of \ControlsLargeFactorRSS{} versus \ControlsLargeDenseRSS{} GiB
(Table~\ref{tab:dense-scaling}). These two-core controls use the same initial
assignments and search limits. They are separate from the four-restart,
32-core two-million-entity public-API experiment below.

\begin{table}[htbp]\centering\small
\begin{tabular}{lrrrrr}
\toprule
Entities & Dense (s) & Factor (s) & Dense GiB & Factor GiB & Cut ratio\\\midrule
4,000 & 1.38 & 1.08 & 0.92 & 0.82 & 1.0003\\
16,000 & 23.79 & 4.52 & 3.93 & 1.36 & 1.0004\\
32,000 & 82.68 & 8.74 & 12.04 & 1.83 & 1.0004\\
\bottomrule\end{tabular}
\caption{Dense versus rank-256 factor optimization on nested ChEMBL samples, using identical initial assignments and search limits on two CPU cores with a 1\% size corridor. Time includes representation construction and optimization; peak process RSS precedes independent scoring. Cut ratio divides the factor solution's original-kernel cut by the dense solution's cut. Entries average three seeds. Both solutions are heuristic.}
\label{tab:dense-scaling}
\end{table}


Equal-time controls complement the iteration-matched results. Their timed
optimizer is stopped externally at the deadline, and only a feasible
checkpoint completed before that deadline is scored. Representation
construction consumes the same budget as search. The supplement also retains
the earlier cooperative-deadline sensitivity cases and their overruns.

Error relative to the optimized cut can be appreciably larger than error
relative to all kernel mass. Reanalysing the 900 existing full-kernel
comparisons, the largest rank-256 cut-relative errors are
\ControlsPlainCutError{}\% for ordinary similarity and
\ControlsPowerCutError{}\% for cubed similarity. These are approximation
errors at tested partitions, not optimization gaps or uniform guarantees.
The supplementary tables retain both normalizations and distinguish the
returned partition from count-preserving perturbations.

Current-version component controls further separate feasible initialization,
balanced assignment, local refinement, restart budget and landmark sampling.
Assignment accounts for most of the observed global reduction: the average
score changes from 1.001 for feasible initialization to 0.626 for assignment
alone, versus 0.625 with four restarts and two refined candidates. Extra search
does not consistently improve nearest-training exposure, because that is not
the ordinary factor objective.
The grouping controls fix and strictly preserve the same groups before
comparing their feasible initialization with ordinary- and cubed-factor
optimization. Their complete global, local and timing results accompany the
rank and dense controls in the supplement.
